% !TEX root = ../master.tex

\section{Hauptteil}
\label{sec:hauptteil}

Dieses Kapitel behandelt die theoretischen Grundlagen und die praktische
Einbindung von \emph{Lezer} als Parser-Generator im Kontext der
Übersetzung von Python-Programmen nach TypeScript. Während in Kapitel~1
die typentheoretischen Grundlagen von TypeScript analysiert wurden und
Kapitel~2 geeignete Datenstrukturen diskutierte, wird nun die syntaktische
Ebene betrachtet: die formale Beschreibung, Analyse und Transformation
von Programmen mittels Grammatiken und Parsern.

\subsection{Einordnung in die Theorie formaler Sprachen}

Die syntaktische Analyse basiert auf dem Modell kontextfreier Grammatiken
(Chomsky-Typ~2). Eine kontextfreie Grammatik ist ein vierfach Tupel
$G = (V, \Sigma, P, S)$ bestehend aus einer Menge von syntaktische Variablen oder Nichtterminalen
$V$, Terminalen oder Token $\Sigma$, Grammatikregeln $P$ und einem Startsymbol $S$
\autocite[S.~171]{hopcroft2006automata}.

Die zentrale Aufgabe eines Parsers besteht darin, für eine Eingabesequenz
ein Ableitungs- bzw.\ Parsebaum-Konstrukt zu bestimmen. Ein Parsebaum
repräsentiert die hierarchische Struktur der Ableitung und ist Grundlage
für weitere semantische Transformationen
\autocite[S.~61f.]{grune2008parsing}.

Für die vorliegende Projektarbeit ist insbesondere die Unterscheidung
zwischen Top-Down- und Bottom-Up-Verfahren relevant.
LR-Verfahren (Left-to-right, Rightmost derivation in reverse)
gehören zur Klasse der deterministischen Bottom-Up-Parser und besitzen
für kontextfreie Grammatiken eine hohe Ausdrucksstärke
\autocite[S.~82f.]{grune2008parsing}.

\subsection{Lezer: Konzept und Architektur}

Lezer ist ein moderner Parser-Generator, der im Umfeld des
CodeMirror-Editors entwickelt wurde. Er erzeugt effiziente
LR-basierte Parser in JavaScript bzw.\ TypeScript
\autocite{lezer_docs}.

Im Unterschied zu klassischen Parsergeneratoren wie Yacc oder Bison
\autocite[S.~287f.]{aho2006compilers} verfolgt Lezer einen stark
inkrementellen Ansatz. Das bedeutet, dass bei Änderungen am Quelltext
nicht der gesamte Parsebaum neu konstruiert wird, sondern nur die
betroffenen Teilbereiche aktualisiert werden.

Formal entspricht die von Lezer verarbeitete Grammatik einer
kontextfreien Grammatik mit expliziten Präzedenz- und
Assoziativitätsangaben. Diese Technik dient der Auflösung
klassischer Mehrdeutigkeiten, etwa bei arithmetischen Ausdrücken
\autocite[S.~210]{hopcroft2006automata}.

\subsection{Grammatikdefinition in Lezer}

Eine Lezer-Grammatik wird deklarativ beschrieben. Beispielhaft sei
eine stark vereinfachte Ausdrucksgrammatik angegeben:

\begin{minted}{typescript}
@top Program { Expression }

Expression {
  Expression "+" Term |
  Term
}

Term {
  Number |
  "(" Expression ")"
}

Number { $[0-9]+ }
\end{minted}

Formal entspricht dies der kontextfreien Grammatik:

\[
\begin{aligned}
Program &\rightarrow Expression \\
Expression &\rightarrow Expression + Term \mid Term \\
Term &\rightarrow Number \mid (Expression)
\end{aligned}
\]

Die linke Rekursion wird von LR-basierten Verfahren direkt unterstützt,
im Gegensatz zu LL-Verfahren, bei denen sie zuvor eliminiert werden
müsste \autocite[S.~212]{aho2006compilers}.

\subsection{Parsebäume und abstrakte Syntaxbäume}

Ein wesentlicher Schritt bei der Übersetzung von Python nach TypeScript
besteht in der Konstruktion eines abstrakten Syntaxbaums (AST).
Der konkrete Parsebaum enthält syntaktische Details, während der AST
auf die semantisch relevanten Knoten reduziert wird
\autocite[S.~358f.]{aho2006compilers}.

Lezer erzeugt intern eine kompakte Baumrepräsentation, die
inkrementell aktualisierbar ist. Für Transformationszwecke wird
dieser Baum traversiert und in eine eigene AST-Struktur überführt,
welche anschließend zur Codegenerierung genutzt wird.

Die formale Korrektheit dieses Vorgehens lässt sich als
Homomorphismus zwischen zwei Baumstrukturen interpretieren:
\[
h : T_{\text{Python}} \rightarrow T_{\text{TypeScript}}
\]
wobei $T_{\text{Python}}$ der durch die Python-Grammatik erzeugte
Baum ist und $T_{\text{TypeScript}}$ die Zielstruktur repräsentiert.

\subsection{Inkrementelles Parsing und Komplexität}

Klassische LR-Parser besitzen im deterministischen Fall lineare
Laufzeit in der Länge der Eingabe
\autocite[S.~73]{hopcroft2006automata}.
Lezer erweitert dieses Modell um inkrementelle Aktualisierung.

Sei $n$ die Länge des Quelltexts und $k$ die Länge einer lokalen
Änderung, so wird im Idealfall nur ein Teilbaum der Größe
$\mathcal{O}(k)$ neu berechnet.
Dies ist insbesondere für editorintegrierte Anwendungen
relevant.

\subsection{Bedeutung für die Python--TypeScript-Übersetzung}

Im Rahmen der Projektarbeit übernimmt Lezer folgende Aufgaben:

\begin{itemize}
    \item Formale Spezifikation der Python-Teilgrammatik
    \item Konstruktion eines syntaktisch korrekten Parsebaums
    \item Bereitstellung einer strukturierten Baumrepräsentation
          für die Transformation
\end{itemize}

Die syntaktische Korrektheit ist Voraussetzung für die spätere
Typtransformation und Codegenerierung. Fehlerhafte oder
mehrdeutige Grammatiken würden direkt zu fehlerhaften
Übersetzungen führen.

\subsection{Zusammenfassung}

Lezer stellt einen modernen, LR-basierten Parser-Generator dar,
der sich aufgrund seiner inkrementellen Architektur besonders
für editornahe Anwendungen eignet. Formal basiert er auf
kontextfreien Grammatiken und deterministischen
Bottom-Up-Verfahren.

Für die vorliegende Studienarbeit bildet Lezer die Brücke
zwischen der formalen Beschreibung der Python-Syntax
und der strukturierten Transformation in typisierte
TypeScript-Programme.

Im folgenden Kapitel wird die konkrete Modellierung der
Python-Grammatik und deren Anpassung an die Anforderungen
der Übersetzung detailliert analysiert.